% sec-09-additional-information.tex - Additional Information (final IND)
\section{Additional Information}

This section addresses the categories of additional information enumerated for an
original Investigational New Drug (IND) submission under 21 CFR \S 312.23(a)(10),
namely drug dependence and abuse potential, radioactive drugs, pediatric studies,
and any other information bearing on the safe and ethical conduct of the proposed
clinical investigation. Because this is a combined IND and Investigational Device
Exemption (IDE) first-in-human study of orally administered daraxonrasib
(RMC-6236) delivered alongside an on-premises large language model (LLM) directed
eight-arm robotic pancreaticoduodenectomy, the Other Information subsection
additionally consolidates the Physical AI governance, cybersecurity, record
retention, and three-tier oversight content that the autonomy-bearing
investigational device requires for Phase 1 acceptance. Each statement below is
grounded in the protocol (v1.0.0, DOI 10.5281/zenodo.20780121), is internally
consistent with the General Investigational Plan, the Proposed Clinical Research,
and the Previous Human Experience sections of this IND (repository v4.3.0, IND
v1.0), and is aligned with FDA Phase 1 guidance \cite{phase1guidance} and with the
regulatory adaptations developed for Physical AI investigations
\cite{cfr312paiuni,cfr050paiuni,ichpaistduni}.

The four enumerated categories of 21 CFR \S 312.23(a)(10) are treated in turn, and
each is given an explicit applicability determination so that the reviewing
division can locate, on a single pass, the basis on which the category is either
satisfied or found not applicable. Drug dependence and abuse potential is addressed
and found to require no abuse-liability program. Radioactive drugs is found not
applicable because the investigational agent contains no radionuclide. Pediatric
studies is addressed by a documented deferral consistent with an adult-only,
serious-and-life-threatening oncology indication that does not arise materially in
children. The Other Information subsection then carries the Physical AI overlay,
which is the principal substantive addition that distinguishes this submission from
a conventional small-molecule Phase 1 IND and which the autonomy-bearing
investigational device makes necessary for safe and ethical conduct. The structure
mirrors the enumeration of the regulation so that completeness can be confirmed by
inspection.

\subsection{Drug Dependence and Abuse Potential}

Daraxonrasib has no known dependence or abuse potential and is not a controlled
substance under the Controlled Substances Act. The investigational drug is an oral,
once-daily RAS(ON) multi-selective (pan-RAS) small-molecule inhibitor that binds
the guanosine triphosphate (GTP) bound active form of mutant RAS in a complex with
cyclophilin A and thereby suppresses downstream effector signaling
\cite{daraxwhipple,onpremwhippl}. Its pharmacological mechanism is targeted
oncogenic-signaling inhibition; it has no activity at the central nervous system
receptors, transporters, or ion channels that mediate reinforcement, euphoria,
sedation, or physical dependence, and it does not belong to a pharmacological class
associated with recreational use, tolerance with escalating self-administration, or
a withdrawal syndrome. The drug carries an FDA Breakthrough Therapy designation
(June 2025) for previously treated metastatic pancreatic ductal adenocarcinoma
(PDAC) with KRAS G12 mutations, and the prior and ongoing human experience
summarized in the Previous Human Experience section has not identified any signal
of abuse, dependence, drug-seeking behavior, or diversion.

The applicability analysis follows the FDA assessment framework for abuse potential.
That framework directs a sponsor to consider, first, the pharmacological class and
mechanism of the molecule and its structural relationship to known drugs of abuse;
second, the receptor-binding and central-nervous-system activity profile; third,
the nonclinical behavioral pharmacology, including any signal in self-administration,
drug-discrimination, and physical-dependence and withdrawal models; and fourth, the
clinical experience to date, including any reports of euphoria, dissociation, mood
elevation, or aberrant drug-seeking behavior. Applied to daraxonrasib, each of these
considerations is negative. The molecule is a RAS(ON) inhibitor with no structural
or pharmacological kinship to opioids, stimulants, cannabinoids, sedative-hypnotics,
dissociatives, or hallucinogens; its molecular target is the active conformation of
mutant RAS rather than any neurotransmitter system implicated in reward; and the
accumulated human experience reflects a tolerability profile dominated by
mechanism-related and gastrointestinal events rather than by any central effect that
would support reinforcement. The conclusion that no abuse-liability program is
required follows directly from the convergence of these four lines of evidence, not
from the absence of any single study.

Accordingly, no abuse-liability program is proposed for this IND. No dedicated
nonclinical abuse-potential studies (for example, self-administration, drug
discrimination, or physical-dependence and withdrawal studies under the FDA
abuse-potential assessment framework) are warranted, and no controlled-substance
scheduling recommendation applies. The investigational product is dispensed under
the accountability controls described in the Chemistry, Manufacturing, and Control
Information and Proposed Clinical Research sections, with drug accountability logs
maintained at the single academic site, so that the limited quantities used for the
planned cohorts (up to eighteen treated participants across three dose levels) are
reconciled from receipt through administration, return, and destruction. Should any
unanticipated central-nervous-system finding suggestive of abuse potential emerge
during the conduct of the study, it would be captured through the routine adverse
event and serious adverse event reporting pathways under 21 CFR \S 312.32 and
reviewed by the Data Safety Monitoring Board (DSMB).

A worked illustration of the accountability arithmetic makes the control concrete.
With a maximum of eighteen treated participants and a once-daily oral regimen, the
study exposure footprint is small and fully reconcilable. The drug accountability
log records, per participant and per cohort, the quantity received, the quantity
dispensed at each visit, the quantity returned unused, and the quantity destroyed,
with each transaction time-stamped and counter-signed. Because the three dose levels
(160, 220, and 300 mg once daily) are dispensed in defined quantities tied to the
twenty-eight-day dose-limiting-toxicity (DLT) window and to the perioperative
pause-and-restart advisory, the expected dispensed quantity for any participant is
computable in advance, and any discrepancy between expected and reconciled quantity
is an audit signal that is reviewed before the next dispensing. This closed-loop
reconciliation, rather than scheduling or a dedicated abuse-deterrent formulation,
is the proportionate control for a non-scheduled targeted oncology agent.

For completeness with respect to 21 CFR \S 312.23(a)(10)(i), the IND states
affirmatively that the investigational agent is not listed in any schedule of the
Controlled Substances Act, that no application for scheduling is pending or
contemplated, and that no labeling, packaging, or distribution control beyond the
ordinary investigational-product accountability of 21 CFR \S 312.57 and \S 312.62 is
required on the basis of abuse potential. The perioperative pause-and-restart of the
drug, in which dosing is held before surgery and resumed on a human-confirmed
advisory, further limits continuous exposure and is itself recorded, so that the
total quantity dispensed, held, resumed, and returned is auditable at every step. No
take-home quantity is dispensed beyond the defined per-visit allotment, and any drug
returned at a visit is reconciled before the next allotment is released. Because the
study is conducted at a single academic site under direct institutional pharmacy
control rather than through outpatient community dispensing, the diversion surface is
small and the accountability chain is short, which is the appropriate level of
control for a non-scheduled targeted oncology agent administered to a closely
monitored surgical population.

\needspace{10\baselineskip}
\begin{table}[htbp]
\centering
\footnotesize
\caption{Abuse-potential assessment summary for daraxonrasib against the four
considerations of the FDA abuse-potential framework, with the determination for
each. The convergence of the four negative determinations supports the conclusion
that no abuse-liability program is required for this IND.}
\label{tab:abuse-potential}
\begin{tabularx}{\textwidth}{L{4.2cm} Y L{2.6cm}}
\toprule
\textbf{Consideration} & \textbf{Finding for daraxonrasib} & \textbf{Determination} \\
\midrule
Pharmacological class and structure & RAS(ON) multi-selective (pan-RAS) small-molecule inhibitor; no structural or class kinship to opioids, stimulants, cannabinoids, sedative-hypnotics, dissociatives, or hallucinogens. & No abuse class \\
Central-nervous-system activity & No activity at receptors, transporters, or ion channels that mediate reinforcement, euphoria, sedation, or physical dependence; target is GTP-bound active mutant RAS in complex with cyclophilin A. & No reward signal \\
Nonclinical behavioral pharmacology & No signal warranting self-administration, drug-discrimination, or physical-dependence and withdrawal studies; no dedicated abuse-potential studies indicated. & Studies not warranted \\
Clinical experience to date & Prior and ongoing human experience (Breakthrough Therapy, June 2025) shows no abuse, dependence, drug-seeking behavior, or diversion. & No clinical signal \\
\bottomrule
\end{tabularx}
\end{table}

\figcaption{Table 10.1. Abuse-potential assessment summary. Each of the four
considerations of the FDA abuse-potential framework returns a negative
determination, supporting the conclusion that no abuse-liability program, no
dedicated nonclinical abuse-potential studies, and no controlled-substance
scheduling recommendation are required for this IND.}

\subsection{Radioactive Drugs}

This subsection is not applicable. Daraxonrasib is not a radioactive drug and
contains no radionuclide. The investigational product is a non-radioactive oral
small-molecule tablet administered once daily; it is neither a positron-emitting
nor a single-photon-emitting radiopharmaceutical, and it is not the subject of a
Radioactive Drug Research Committee (RDRC) determination under 21 CFR \S 361.1. No
radiation dosimetry, no radiation safety committee review, and no radioactive
drug-specific manufacturing, handling, storage, or waste-disposal controls are
required for the investigational agent. No information under 21 CFR
\S 312.23(a)(10)(ii) is therefore required.

Imaging performed during the study is part of routine clinical care and is not an
investigational radioactive drug administered under this IND. Diagnostic computed
tomography, magnetic resonance imaging, and any standard-of-care nuclear medicine
staging that may precede enrollment are performed and interpreted according to
institutional standards and are outside the scope of this subsection; any
radiopharmaceutical used in such standard-of-care imaging is administered under the
institution's own authorizations and not under this IND. To avoid ambiguity for the
reviewing division, the operative and device elements of the study are also
addressed. The intraoperative imaging, hemostasis, and drain placement of operative
phase P8 are optical and instrumental and involve no administered radioactivity; the
eight-arm robotic system and its 640 sensor channels are electromechanical and
optical instruments, not radiation-emitting devices within the meaning of a
radioactive drug determination. The vascular no-fly gate, the force-envelope
controls, and the cyber-physical telemetry stream described in the Other Information
subsection are likewise non-radioactive. No element of the investigational
product or the investigational device introduces a radionuclide into the
participant, and the radioactive-drug category is therefore satisfied as not
applicable in its entirety.

\subsection{Pediatric Studies}

The study population is restricted to adults aged eighteen years or older. Eligible
participants have histologically or cytologically confirmed PDAC that is resectable
or borderline-resectable, a documented KRAS G12 mutation (G12D, G12V, or G12R), an
Eastern Cooperative Oncology Group (ECOG) performance status of 0 to 1, and are
candidates for pancreaticoduodenectomy (Whipple resection), consistent with the
broad pan-KRAS eligibility framework of the RASolute 302 program
\cite{daraxwhipple}. No participant under the age of eighteen will be enrolled,
consented, screened, dosed, or operated upon under this protocol, and pediatric
assent provisions are therefore inapplicable in practice even though the
informed-consent architecture (21 CFR part 50 with ICH E6(R3)) provides for assent
where it would otherwise apply \cite{cfr050paiuni}.

A first-in-human Phase 1 evaluation in a pediatric population is neither warranted
nor proposed at this stage, and a pediatric study plan is deferred. The deferral
rationale rests on three considerations. First, the disease under study, PDAC, is
overwhelmingly a disease of older adults; pancreatic ductal adenocarcinoma in
children is exceedingly rare, the five-year overall survival of PDAC is below
thirteen percent, and the condition is the third leading cause of cancer death in
the United States and the fourth in the European Union \cite{pdac060s2030}, so the
target indication does not arise materially in the pediatric population. Second,
this is a first-in-human study of an autonomy-bearing investigational device
combined with an investigational dosing regimen; the prudent and ethical sequence
is to establish the device and procedure serious adverse event profile, the maximum
tolerated dose (MTD) and recommended Phase 2 dose (RP2D), and the technical
feasibility of the LLM-directed robotic Whipple in adults before any consideration
of younger populations, and full autonomy is in any case prohibited at this stage
under 21 CFR \S 312.21(e) \cite{cfr312paiuni}. Third, the staged-autonomy model of
this trial class confines the present study to clinician-controlled teleoperation
and supervised operation (Stage 1), reserving higher-autonomy operation for future
study; extension to a pediatric population, if ever pursued, would follow successful
completion of the adult program and would be the subject of a separate pediatric
study plan with its own age-appropriate formulation, dosing, device sizing, and
assent and parental-permission provisions.

The deferral is also coherent with the device-engineering constraints of the
investigational system. The eight-arm robotic platform, with its 56 degrees of
freedom and its fixed force-envelope caps of no more than 3 N per arm and no more
than 18 N cumulative, is validated for the anatomy of an adult pancreaticoduodenal
field and for the vessel-proximity tolerances of an adult superior mesenteric vein,
portal vein, hepatic artery, celiac artery, and superior mesenteric artery. The
vascular no-fly margins of 1.0 mm at the superior mesenteric vein and portal vein
and 1.5 mm at the hepatic artery, celiac artery, and superior mesenteric artery,
together with the Phase 0 simulation-validation gate, are calibrated to adult
geometry. A pediatric extension would require revalidation of the entire
cyber-physical envelope, the digital twin, and the simulation corpus against
pediatric anatomy and would not be a simple eligibility amendment. This engineering
reality reinforces the clinical and ethical case for deferring any pediatric
evaluation until the adult program has established the safety and feasibility
foundation on which any such revalidation would have to build.

Because this is an investigation of a serious and life-threatening oncology
indication that does not occur meaningfully in children, the study qualifies for
the conditions under which a pediatric assessment is not required for an adult-only
first-in-human Phase 1 program; any subsequent pediatric commitments would be
addressed at the appropriate later stage of development in consultation with the
FDA. No pediatric data are submitted with this original IND.

\subsection{Other Information}

The autonomy-bearing investigational device requires governance, cybersecurity,
record-retention, and oversight controls beyond those of a conventional
small-molecule Phase 1 study. This subsection consolidates that Physical AI overlay.
The on-premises LLM operates strictly as a second-opinion advisory oracle: it is
hash-pinned to a deterministic seed and a repository commit, isolated outside the
vendor kinematic stack, and never granted control authority over the robotic arms;
full autonomy is prohibited under 21 CFR \S 312.21(e) \cite{cfr312paiuni}, and every
proposed robot-patient command passes a verification-before-generation ten-gate
suite returning accept, block, or escalate, aligned with the Verification Before
Generation in Physical AI Oncology Trials Act of 2026 (H.R. 9510)
\cite{congress9510}. The governance design follows the trust and accountability
principles developed for clinical Physical AI \cite{clintrustpai,paiautospons} and
the site-readiness and national-platform work for Physical AI oncology trials
\cite{natlpaiplatf,paisitedocpk}.

\subsubsection{Physical AI Governance and the Three-Tier Oversight Structure}

Independent safety oversight operates on three coordinated tiers layered over the
reviewing Institutional Review Board (IRB) and the Sponsor-Investigator. The DSMB,
composed of members independent of the conduct of the study with surgical,
oncologic, biostatistical, and device-safety expertise, reviews accumulating safety
data at each cohort boundary and at any triggered review and holds binding halt
authority. The Independent Safety Monitor (ISM) is designated for each procedure and
provides real-time, case-level safety oversight independent of the operating team,
with the authority to call an immediate stop. The Physical AI Safety Review
Committee meets on a cadence of ninety days or less (and ad hoc on any triggering
event) to review the autonomy-bearing behavior of the system, including the
telemetry and audit record, the verification-before-generation gate-surface
decisions, the Unified Safety Level (USL) readiness rating, and any software change,
and recommends USL reassessment or oversight changes to the Sponsor-Investigator.
The three tiers are coordinated so that no autonomy-bearing change and no escalation
decision proceeds without appropriate independent review. The Sponsor-Investigator
(ChemicalQDevice, Kevin Kawchak, Chief Executive Officer) holds the combined IND and
IDE; the conflict of interest arising from the Sponsor-Investigator's relationship
to the developer of the Physical AI system is disclosed in full under 21 CFR part 54
and is mitigated structurally by this independent oversight, so that safety and
escalation decisions are not made by a financially interested party alone.

The coordination among the tiers is deliberate and is designed to remove any single
point of discretion over a safety-critical decision. The reviewing IRB exercises
human-subjects protection and reviews the protocol, the informed consent, and the
significant-risk device determination, and it may suspend or terminate its approval.
The DSMB applies the prespecified halt rules at each cohort boundary and at any
triggered review and recommends continue, modify, pause, or halt. The ISM exercises
real-time authority at the level of the individual procedure and may stop a case
without waiting for a scheduled review. The Physical AI Safety Review Committee
exercises the autonomy-specific oversight that a conventional DSMB is not equipped to
perform, reading the telemetry and audit record and the gate-surface decisions and
recommending USL reassessment when the evidence requires it. Because the
Sponsor-Investigator and the developer relationship are the same financial interest,
the structural mitigation is that the binding halt authority, the per-procedure stop
authority, and the autonomy-readiness reassessment authority all reside in bodies
independent of that interest, so that no safety-critical determination can be made
by the interested party alone. Table~\ref{tab:oversight-bodies} summarizes the
oversight bodies, their cadence, and their authority, and Figure~20 renders the
governance structure.

\needspace{12\baselineskip}
\begin{table}[htbp]
\centering
\footnotesize
\caption{Three-tier independent safety oversight and the human-subjects-protection
structure for the combined IND and IDE study. Each body's review cadence and the
authority it exercises are stated; the binding halt rules drive the continue,
modify, pause, or halt decision.}
\label{tab:oversight-bodies}
\begin{tabularx}{\textwidth}{L{3.4cm} L{3.2cm} Y}
\toprule
\textbf{Body} & \textbf{Review cadence} & \textbf{Role and authority} \\
\midrule
Sponsor-Investigator (ChemicalQDevice, K. Kawchak) & Continuous & Holds combined IND and IDE; protocol custody, safety reporting under 21 CFR parts 312 and 812, and compliance assurance; conflict of interest disclosed and managed under 21 CFR part 54. \\
Reviewing IRB & Per submission and continuing review & Human-subjects protection; reviews protocol, consent, and the significant-risk device determination; may suspend or terminate approval. \\
Data Safety Monitoring Board (DSMB) & Each cohort boundary and any triggered review & Independent surgical, oncologic, biostatistical, and device-safety members; applies the prespecified halt rules; recommends continue, modify, pause, or halt; binding halt authority. \\
Independent Safety Monitor (ISM) & Per procedure, real-time & Case-level safety oversight independent of the operating team; authority to call an immediate stop. \\
Physical AI Safety Review Committee & Cadence of 90 days or less, plus ad hoc & Reviews telemetry and audit records, verification-before-generation gate-surface decisions, USL readiness, and software changes; recommends USL reassessment or oversight changes. \\
\bottomrule
\end{tabularx}
\end{table}

\figcaption{Table 10.2. The oversight bodies, their cadence, and their authority for
the combined IND and IDE first-in-human study.}

\needspace{8\baselineskip}
\begin{mermaidfig}
\node[mmgoal] (si)    at (4.6,0)      {Sponsor-Investigator\\ChemicalQDevice (K. Kawchak)\\COI disclosed, 21 CFR part 54};
\node[mmproc] (irb)   at (-0.4,-3.0)  {Reviewing IRB\\(human-subjects protection)};
\node[mmproc] (dsmb)  at (4.6,-3.0)   {DSMB\\cohort-boundary + triggered;\\halt authority};
\node[mmproc] (ism)   at (9.6,-3.0)   {Independent Safety Monitor\\per-procedure, real-time;\\stop authority};
\node[mmproc] (pais)  at (-0.4,-6.0)  {Physical AI Safety Review\\Committee ($\leq$ 90-day cadence)};
\node[mmdark] (halt)  at (4.6,-6.0)   {Binding halt rules:\\device SAE $\geq$ 1 of 3 in a cohort\\OR 2 device-related deaths};
\node[mmdec]  (dec)   at (4.6,-9.2)   {Continue, modify,\\pause, or halt};
\draw[mmedge] (si.south) to[out=-90,in=70,looseness=0.6] (irb.north);
\draw[mmedge] (si) -- (dsmb);
\draw[mmedge] (si.south) to[out=-90,in=110,looseness=0.6] (ism.north);
\draw[mmedge] (si.west) to[out=180,in=90,looseness=0.6] (pais.north);
\draw[mmedge] (dsmb) -- (halt);
\draw[mmedge] (halt) -- (dec);
\draw[mmedge] (ism.south) to[out=-90,in=20,looseness=0.6] (dec.north east);
\draw[mmedge] (pais.south) to[out=-90,in=200,looseness=0.6] (dec.north west);
\end{mermaidfig}
\figcaption{Figure 20. Study governance and three-tier safety oversight. The
Sponsor-Investigator (ChemicalQDevice, Kevin Kawchak, with the conflict of interest
disclosed under 21 CFR part 54) operates under a reviewing IRB and three independent
bodies: the Data Safety Monitoring Board (cohort-boundary and triggered reviews,
with halt authority), the Independent Safety Monitor (per-procedure, real-time, with
stop authority), and the Physical AI Safety Review Committee (a cadence of 90 days or
less). The binding halt rules, a device-related serious adverse event in at least
one of three within a cohort or two device-related deaths at any point, drive the
continue, modify, pause, or halt decision.}

\subsubsection{Binding Halt Rules and Stopping Authority}

Two prespecified, binding halt rules govern the study and are applied by the DSMB.
First, a device-related or procedure-related serious adverse event in at least one
of three participants within a dose cohort triggers a mandatory DSMB review with
authority to halt. Second, two device-related deaths at any point in the study
trigger a mandatory DSMB review with authority to halt. On invocation of either
rule, enrollment and any patient-facing robotic activity stop pending review. These
rules operate independently of, and in addition to, the dose-finding rules of the
3+3 escalation (dose levels 160, 220, and 300 mg once daily, with a twenty-eight-day
dose-limiting-toxicity window) and the per-procedure stop authority of the
Independent Safety Monitor. The cyber-physical safety envelope provides further
hard stops that are not discretionary: force caps of no more than 3 N per arm and no
more than 18 N cumulative across the eight arms, a 10 kHz heartbeat with a
sixty-four-byte frame and a one-hundred-microsecond watchdog, arm-park within fifty
microseconds, a cross-arm emergency stop within three milliseconds to a commanded
0.0 N, and a hardware emergency stop within five hundred milliseconds consistent
with 21 CFR \S 312.404. The vascular no-fly gate samples at 10 kHz and enforces a
soft-warning advisory zone at 3.0 mm, blocking no-fly margins of 1.0 mm at the
superior mesenteric vein and portal vein and 1.5 mm at the hepatic artery, celiac
artery, and superior mesenteric artery, and a hard-stop emergency-stop zone at 5.0
mm. Any breach of a hard cyber-physical limit, a failure of the
verification-before-generation gate surface, or a loss of audit integrity
constitutes a clinical-hold ground for the autonomy-bearing system, on which all
patient-facing robotic activity ceases and the decommissioning and data-preservation
procedure is executed.

The interaction between the binding clinical halt rules and the deterministic
cyber-physical limits is worth making explicit, because the two operate on different
time scales and through different mechanisms. The cyber-physical limits are
microsecond-to-millisecond hardware and firmware interlocks that act before any human
or committee can deliberate; the cross-arm emergency stop reaches a commanded 0.0 N
within three milliseconds and the hardware emergency stop completes within five
hundred milliseconds. These are not discretionary and do not wait for a vote. The
binding halt rules, by contrast, are deliberative thresholds that the DSMB applies
after the fact to accumulating safety data, and they determine whether the study as a
whole continues, is modified, is paused, or is halted. A single hard-stop activation
in the operating field is handled by the envelope in milliseconds and is recorded in
the safety-limit-event record; the pattern of such activations, together with the
clinical outcomes, is what the DSMB reviews against the binding halt rules at the
cohort boundary. The two mechanisms are complementary: the envelope protects the
individual participant in real time, and the halt rules protect the cohort and the
study by governing whether exposure continues. Table~\ref{tab:halt-envelope} sets out
the deterministic envelope thresholds, their latency budgets, and the regulatory
anchor for each.

\needspace{12\baselineskip}
\begin{table}[htbp]
\centering
\footnotesize
\caption{Deterministic cyber-physical safety envelope thresholds and latency
budgets, distinguished from the deliberative binding halt rules. The envelope acts
in real time at the level of the individual procedure; the binding halt rules are
applied by the DSMB to accumulating data at the cohort and study level.}
\label{tab:halt-envelope}
\begin{tabularx}{\textwidth}{L{4.6cm} L{4.2cm} Y}
\toprule
\textbf{Control} & \textbf{Threshold or latency} & \textbf{Function and anchor} \\
\midrule
Per-arm force cap & $\leq$ 3 N per arm & Limits tissue-contact force; recorded in the safety-limit-event record. \\
Cumulative force cap & $\leq$ 18 N across eight arms & Bounds aggregate force across 56 degrees of freedom. \\
Control heartbeat & 10 kHz, 64-byte frame, 100 us watchdog & Liveness of the cyber-physical loop; loss triggers arm-park. \\
Arm-park latency & $\leq$ 50 us & Safe halt of an arm on watchdog timeout. \\
Cross-arm emergency stop & $\leq$ 3 ms to commanded 0.0 N & Coordinated stop of all eight arms. \\
Hardware emergency stop & $\leq$ 500 ms & Independent hardware-level stop (21 CFR \S 312.404). \\
Vascular no-fly gate (soft) & 3.0 mm advisory, 10 kHz sampling & Soft-warning advisory zone at the five named vessels. \\
Vascular no-fly gate (block) & 1.0 mm SMV and PV; 1.5 mm HA, CA, SMA & No-fly block margin at the five named vessels. \\
Vascular no-fly gate (hard-stop) & 5.0 mm emergency-stop zone & Triggers emergency stop on incursion. \\
\bottomrule
\end{tabularx}
\end{table}

\figcaption{Table 10.3. The deterministic cyber-physical safety envelope. These
real-time interlocks act at the individual-procedure level on microsecond-to-
millisecond time scales and are not discretionary; the DSMB binding halt rules act
on accumulating data at the cohort and study level.}

A representative 100-tick sample of the vascular no-fly gate during dissection,
drawn from the protocol, illustrates how the gate verdicts distribute in practice:
of one hundred sampled ticks, eighty-three return clear, six return a soft-warning
advisory, six return a no-fly block, and five return a hard-stop emergency-stop
verdict. The eleven blocking or stopping verdicts in this sample are each written to
the safety-limit-event record with the vessel, the margin, and the latency of the
response, so that the DSMB can review both the rate of gate engagement and the
fidelity of the response when it convenes at the cohort boundary. This is the
mechanism by which a real-time interlock event becomes part of the deliberative
record that governs the binding halt determination.

\subsubsection{Cybersecurity and System Integrity}

Cybersecurity is treated as a patient-safety control rather than an administrative
adjunct. The LLM-directed system is deployed on premises and isolated outside the
vendor kinematic stack, with no requirement for external network connectivity during
a procedure; the advisory model is hash-pinned to a deterministic seed and a
repository commit so that the validated configuration in use is exactly the
configuration that was tested and recorded. Access to identifiable records and to
the cyber-physical control surface is governed by a deny-by-default authorization
model, under which no role can read, write, or export a record except where an
access right has been explicitly granted, and electronic signatures are bound to the
records they authenticate. Every access, change, and export is written to a
hash-chained audit trail in which each entry cryptographically incorporates the hash
of the prior entry, so that the record is tamper-evident and any alteration is
detectable. A cybersecurity incident is one of the six Physical AI reporting
triggers under 21 CFR \S 312.32(g) and is reportable within fifteen calendar days,
or within seven calendar days where there is potential for patient harm; the audit
trail is preserved from twenty-four hours before to seventy-two hours after each
Physical AI adverse event under 21 CFR part 11.

The six Physical AI reporting triggers under 21 CFR \S 312.32(g) together extend the
conventional Subpart J reporting framework to the autonomy-bearing behaviors of the
system, and each carries a defined clock so that the obligation is unambiguous. A
serious Physical AI adverse event is reported on the fatal-or-life-threatening
seven-calendar-day clock or the other-serious-and-unexpected fifteen-calendar-day
clock, parallel to the conventional suspected-adverse-reaction reporting of 21 CFR
\S 312.32. A system-drug interaction, in which the device behavior and the
investigational drug interact in a manner bearing on safety, is reported within
fifteen calendar days. A cybersecurity incident is reported within fifteen calendar
days, or within seven where there is potential for patient harm. Sustained AI model
degradation, defined as degradation persisting across at least three consecutive
procedures or twenty-four cumulative hours, is reported because it indicates a drift
of the advisory oracle away from its validated behavior. A sim-to-real divergence
exceeding twice the tolerance, namely a trajectory error above 4 mm or a force error
above 1.0 N, is reported because it indicates that the digital twin and the physical
system have parted company beyond the validated envelope. A digital-twin discrepancy
is reported on the same pathway. The common feature of the six triggers is that each
denotes a departure of the autonomy-bearing system from its validated, deterministic
behavior, whether that departure is adversarial (the cybersecurity incident), gradual
(the sustained model degradation), or a loss of fidelity between the model and the
physical system (the sim-to-real divergence and the digital-twin discrepancy). The
seven-and-fifteen-day clocks parallel the conventional 21 CFR \S 312.32 reporting so
that the autonomy overlay does not create a separate or slower pathway, and the
seven-day cybersecurity clock where there is potential for patient harm reflects the
elevated urgency of a control-surface compromise. Table~\ref{tab:pai-triggers}
enumerates the six triggers, their definitions, and their clocks.

\needspace{12\baselineskip}
\begin{table}[htbp]
\centering
\footnotesize
\caption{The six Physical AI reporting triggers under 21 CFR \S 312.32(g), with the
definition and the reporting clock for each. These triggers extend the conventional
Subpart J reporting framework to the autonomy-bearing behaviors of the LLM-directed
robotic system.}
\label{tab:pai-triggers}
\begin{tabularx}{\textwidth}{L{4.0cm} Y L{2.8cm}}
\toprule
\textbf{Trigger} & \textbf{Definition} & \textbf{Clock} \\
\midrule
Serious Physical AI adverse event & A serious adverse event in which the autonomy-bearing system is implicated, on the conventional 21 CFR \S 312.32 seriousness and expectedness logic. & 7 days (fatal or life-threatening); 15 days (other serious, unexpected) \\
System-drug interaction & A safety-relevant interaction between the device behavior and the investigational drug. & 15 days \\
Cybersecurity incident & A compromise or attempted compromise of the system, records, or control surface. & 15 days; 7 days if potential for patient harm \\
Sustained AI model degradation & Degradation persisting across at least three consecutive procedures or 24 cumulative hours. & 15 days \\
Sim-to-real divergence beyond 2x tolerance & Trajectory error above 4 mm or tip-force error above 1.0 N. & 15 days \\
Digital-twin discrepancy & A discrepancy between the digital twin and the physical system bearing on safety. & 15 days \\
\bottomrule
\end{tabularx}
\end{table}

\figcaption{Table 10.4. The six Physical AI reporting triggers under 21 CFR
\S 312.32(g). Each trigger carries a defined reporting clock, and each generates an
entry in the hash-chained audit trail preserved from 24 hours before to 72 hours
after the event under 21 CFR part 11.}

\subsubsection{The Seven Physical AI Record Types and 21 CFR Part 11}

Seven Physical AI record types are captured and retained for this trial class, in
addition to the conventional clinical-trial records. Table~\ref{tab:record-types}
enumerates each record type, its content, and its primary regulatory anchor. All
seven are bound to a deterministic seed and a repository commit by the hash-chained
audit trail (record type seven), so that any reported behavior can be replayed
deterministically and any requested record reconstructed. All electronic records,
signatures, and the cyber-physical telemetry stream are maintained under 21 CFR part
11, with validated systems, secure time-stamped audit trails that do not obscure
prior entries, role-based access controls, and electronic signatures bound to the
records they authenticate. Records are retained for the period required by 21 CFR
\S 312.57, namely at least two years after a marketing application is approved for
the indication or, if no application is filed or approved, two years after the
investigation is discontinued and the FDA is notified; the raw, losslessly
reconstructable command and sensor tier is preserved so that a requested record can
be reconstructed on FDA request under 21 CFR \S 312.57. Data released for analysis,
deposition, or publication are de-identified by the HIPAA Safe Harbor method, with
removal of all eighteen categories of protected health information identifiers.

The volume and tiering of the telemetry record warrant a concrete description,
because the determinism guarantee depends on retaining enough of the stream to
reconstruct any reported behavior. The system carries 640 sensor channels, eighty
per arm across eight arms, sampling force at 100 kHz and other channels at 10 kHz.
The sensor and telemetry record (record type three) is tiered from compact summaries,
suitable for routine review, down to a raw, losslessly reconstructable level that
preserves the full stream around any event of interest. The robot command record
(record type two) captures each command issued to the eight arms and their 56 degrees
of freedom with timestamps sufficient for deterministic replay, and the LLM advisory
record (record type one) captures each prompt, the supplied context, and the advisory
output of the second-opinion oracle, hash-pinned to seed and commit and always
advisory rather than control-bearing under 21 CFR \S 312.21(e). The
verification-before-generation gate-surface record (record type five) captures each
accept, block, or escalate decision of the ten-gate suite; the human-override and
intervention record (record type six) captures each manual override, each takeover by
the backup operator, and each stop invoked by the Independent Safety Monitor; and the
safety-limit-event record (record type four) captures the emergency-stop activations,
the force-envelope excursions, and the vascular no-fly gating events. The
hash-chained audit trail (record type seven) binds the preceding six together so that
the entire record is tamper-evident and deterministically replayable.

\needspace{14\baselineskip}
\begin{table}[htbp]
\centering
\footnotesize
\caption{The seven Physical AI record types captured and retained for the
LLM-directed robotic Whipple, in addition to the conventional clinical-trial
records. All records are retained under 21 CFR part 11 for the period required by 21
CFR \S 312.57.}
\label{tab:record-types}
\begin{tabularx}{\textwidth}{L{0.7cm} L{3.6cm} Y}
\toprule
\textbf{No.} & \textbf{Record type} & \textbf{Content and regulatory anchor} \\
\midrule
1 & LLM advisory record & Each prompt, the supplied context, and the advisory output of the second-opinion oracle, hash-pinned to seed and commit; advisory only, never control authority (21 CFR \S 312.21(e)). \\
2 & Robot command record & Each command issued to the eight arms (56 degrees of freedom), with timestamps for deterministic replay. \\
3 & Sensor and telemetry record & The cyber-physical stream from 640 sensor channels (80 per arm; 100 kHz force, 10 kHz other), tiered from compact summaries to the raw losslessly reconstructable level. \\
4 & Safety-limit-event record & Emergency-stop activations and latencies, positional and force-envelope excursions ($\leq$ 3 N per arm, $\leq$ 18 N cumulative), and vascular no-fly gating events (soft 3.0 mm, no-fly 1.0 to 1.5 mm, hard-stop 5.0 mm). \\
5 & Verification-before-generation gate-surface record & Each accept, block, or escalate decision of the ten-gate suite, aligned with H.R. 9510 (Verification Before Generation in Physical AI Oncology Trials Act of 2026). \\
6 & Human-override and intervention record & Each manual override, takeover by the backup operator, and stop invoked by the Independent Safety Monitor. \\
7 & Hash-chained audit trail & Cryptographically chained binding of records 1 to 6 to a deterministic seed and repository commit; tamper-evident; preserved from $-$24 h to $+$72 h around each Physical AI adverse event (21 CFR part 11). \\
\bottomrule
\end{tabularx}
\end{table}

\figcaption{Table 10.5. The seven Physical AI record types, their content, and their
regulatory anchors. Record type seven binds the preceding six to a deterministic
seed and repository commit so that any reported behavior can be replayed
deterministically.}

\subsubsection{Quality Assurance, Re-Validation, and Decommissioning}

Quality assurance spans the clinical and the cyber-physical domains. Before every
procedure, the pre-procedure safety matrix is verified, the USL readiness rating is
confirmed at or above the deployment threshold, and the emergency-stop, positional,
and force envelopes are checked. The Phase 0 simulation-validation gate, which must
be satisfied before any patient-facing robotic activity, requires a USL of at least
7.0, at least one thousand simulations across at least two independent frameworks
(for example Isaac, MuJoCo, Gazebo, or PyBullet), and a sim-to-real agreement of
less than 2 mm in trajectory and less than 0.5 N in tip force. After any software
update, model change, or configuration change to the LLM-directed system, a Phase 0
re-validation is performed and the USL is reassessed before any further
patient-facing use; no update reaches a participant without re-validation.

The re-validation discipline is the operational counterpart of the hash-pinning
guarantee. Because the advisory model is hash-pinned to a deterministic seed and a
repository commit, any change to the model, the configuration, or the surrounding
software produces a new commit and therefore a new validated configuration that has
not yet passed the Phase 0 gate. The rule that no update reaches a participant
without re-validation closes the loop: a change cannot be deployed to a patient-facing
procedure until the new configuration has satisfied the USL threshold of at least
7.0, the simulation-count threshold of at least one thousand runs across at least two
frameworks, and the sim-to-real agreement thresholds of less than 2 mm in trajectory
and less than 0.5 N in tip force. The Physical AI Safety Review Committee reviews each
such change and the associated USL reassessment on its cadence of ninety days or
less, and ad hoc on any triggering event, so that the autonomy-readiness rating in
force at any time corresponds to a configuration that has actually been validated.

On any clinical hold or closure, enrollment and patient-facing robotic activity stop
immediately, affected participants are transitioned to appropriate clinical care,
and a controlled decommissioning and data-preservation procedure is executed under
which the robotic system is locked out of patient contact, the deterministic
configuration and repository commit are recorded, and the complete hash-chained
command and telemetry record is sealed and retained for the regulatory retention
period. The computational evidence supporting this assurance approach includes prior
in silico and digital-twin work by the author, namely a sixty-second PDAC Whipple
with daraxonrasib simulation, an AI digital-twin PDAC simulation, a quantitative
systems pharmacology metastatic-PDAC simulation with verification, validation, and
uncertainty quantification, a one-hundred-thousand-patient in silico five-arm Phase 3
PDAC study in triplicate, and end-to-end PDAC digital-twin trial proposals
\cite{qspmetpancre,fdadigtwinpc,chatgpt100kp,pdacdigtwinp}, together with a
sixteen-LLM code-generation evaluation against FDA adverse-event reporting data
\cite{llmcomp16fda} and the predictive Physical AI oncology modeling that informs the
sim-to-real tolerances \cite{paipredict4x}.

The perioperative daraxonrasib pause-and-restart advisory is an instructive worked
example of how the advisory oracle, the human confirmation step, and the
record-retention discipline operate together within these quality boundaries. The
drug is paused before surgery, and the restart day is determined from an
LLM-bound, human-confirmed advisory whose inputs are the pancreaticojejunostomy
ISGPS fistula grade (A, B, or C) and the serum daraxonrasib trough relative to a
0.5 ng/mL threshold, with outputs of restart at day seven, day fourteen, or day
twenty-one, or hold. A thirty-two-iteration deterministic sweep recorded in the
protocol returns a restart at day seven in twenty-nine of thirty-two iterations
(90.6 percent, corresponding to ISGPS Grade A with a sub-threshold trough), a restart
at day fourteen in three of thirty-two iterations (9.4 percent, corresponding to
ISGPS Grade B or a delayed profile), and no iteration returning day twenty-one, with
a baseline serum trough of 0.45 ng/mL (sub-threshold) across all iterations. The
advisory is recorded in the LLM advisory record (record type one), the human
confirmation is recorded in the human-override and intervention record (record type
six), and both are bound by the hash-chained audit trail (record type seven), so that
the determinism of the sweep can be reproduced on FDA request and the human-in-the-
loop character of the decision is preserved in the record. This example shows the
advisory oracle functioning exactly as 21 CFR \S 312.21(e) requires: informing, never
controlling, a clinical decision that a human makes and records.

\subsection{Selected References}

The following selected references support the additional information in this section
and are cited in full in the back matter. The FDA Phase 1 guidance frames the
content and graded rigor appropriate to an original first-in-human IND
\cite{phase1guidance}, and the daraxonrasib and on-premises robotic Whipple sources
establish the drug mechanism and the device context \cite{daraxwhipple,onpremwhippl}.
The Physical AI regulatory adaptations of 21 CFR parts 312 and 50 and of ICH
guidance govern the autonomy-bearing overlay, the opt-out, and the Subpart J
reporting and record-retention controls \cite{cfr312paiuni,cfr050paiuni,ichpaistduni}.
The Verification Before Generation in Physical AI Oncology Trials Act of 2026 (H.R.
9510) anchors the verification-before-generation gate surface \cite{congress9510},
and the trust, accountability, and site-readiness literature for clinical Physical AI
informs the governance and oversight design
\cite{clintrustpai,paiautospons,natlpaiplatf,paisitedocpk}. The PDAC epidemiology and
the prior computational evidence supporting the deferral and assurance rationale are
drawn from the indication and digital-twin sources
\cite{pdac060s2030,qspmetpancre,fdadigtwinpc,chatgpt100kp,pdacdigtwinp,llmcomp16fda,paipredict4x}.
A complete reference list appears in the back matter.
